
\subsubsection{SHMEMG\_MY\_PE}\label{subsec:shmemg_my_pe}

Returns the \ac{PE} number of the calling \ac{PE}.

\begin{apidefinition}
\begin{Csynopsis}
SHMEMG_DEVICE int shmemg_my_pe(void);
\end{Csynopsis}

\begin{apiarguments}
None.
\end{apiarguments}

\apidescription{
Callable from within a \ac{GPU} kernel.
The \FUNC{shmemg\_my\_pe} routine returns the \ac{PE} number of the calling
\ac{PE}. The value is an integer between 0 and \texttt{npes - 1}, where
\texttt{npes} is the total number of \acp{PE} in the program. The returned value
is identical to the value returned by \FUNC{shmem\_my\_pe} on the host for the
same \ac{PE}.
}

\apireturnvalues{
Integer --- between 0 and \texttt{npes - 1}.
}

\apinotes{
All threads within a kernel launched by the same \ac{PE} shall observe the same
return value.
}
\end{apidefinition}

\subsubsection{SHMEMG\_N\_PES}\label{subsec:shmemg_n_pes}

Returns the number of \acp{PE} in the \openshmem program.

\begin{apidefinition}
\begin{Csynopsis}
SHMEMG_DEVICE int shmemg_n_pes(void);
\end{Csynopsis}

\begin{apiarguments}
None.
\end{apiarguments}

\apidescription{
Callable from within a \ac{GPU} kernel.
The \FUNC{shmemg\_n\_pes} routine returns the total number of \acp{PE} in the
\openshmem program. The returned value is identical to the value returned by
\FUNC{shmem\_n\_pes} on the host.
}

\apireturnvalues{
Integer --- total number of \acp{PE} in the program.
}

\apinotes{
All threads within a kernel launched by the same \ac{PE} shall observe the same
return value.
}
\end{apidefinition}

\paragraph*{\textbf{Example}}\mbox{}\par\nobreak
The following program launches a device kernel that queries the calling \ac{PE}
number and the total number of \acp{PE} from within the kernel using
\FUNC{shmemg\_my\_pe} and \FUNC{shmemg\_n\_pes}, writing the results to a
symmetric buffer in the \ac{GPU} symmetric heap.

\lstinputlisting[style=SourceListingsDefault]{example_code/gc_setup_example.c}

\subsubsection{SHMEMG\_PTR}\label{subsec:shmemg_ptr}

Returns a pointer to a symmetric data object on a specified \ac{PE} that is
directly accessible from within a \ac{GPU} kernel.

\begin{apidefinition}
\begin{Csynopsis}
SHMEMG_DEVICE void *shmemg_ptr(const void *dest, int pe);
\end{Csynopsis}

\begin{apiarguments}
\apiargument{IN}{dest}{Symmetric address of the remotely accessible data object
to be referenced. Shall reside in the \ac{GPU} symmetric heap.}
\apiargument{IN}{pe}{PE number of the remote PE.}
\end{apiarguments}

\apidescription{
Callable from within a \ac{GPU} kernel.
The \FUNC{shmemg\_ptr} routine returns an address that may be used to directly
reference the symmetric data object \VAR{dest} on \ac{PE} \VAR{pe} from within a
\ac{GPU} kernel. This address can be used by the calling \ac{GPU} thread to
perform ordinary loads and stores to the remote data object, bypassing the
\ac{RMA} routines. This is the \ac{GPU}-centric analogue of \FUNC{shmem\_ptr}
in the \openshmem specification.

The routine returns a valid address only when the symmetric data object on the
remote \ac{PE} \VAR{pe} is directly accessible from the calling \ac{PE}'s
\ac{GPU}, for example when the remote \ac{PE} resides on the same node and its
symmetric heap is mapped into the accessing \ac{GPU}'s address space. When the
remote data object is not directly accessible, the routine returns a null
pointer.

Accesses performed through a pointer returned by \FUNC{shmemg\_ptr} are not
ordered or completed by \FUNC{shmemg\_fence} or \FUNC{shmemg\_quiet} with
respect to \ac{RMA} operations. Ordering and visibility of such direct accesses
are governed by the memory model of the underlying \ac{GPU} and the memory
ordering routines described in Section~\ref{subsec:gc_ordering}.
}

\apireturnvalues{
A pointer to the symmetric data object \VAR{dest} on \ac{PE} \VAR{pe} when that
object is directly accessible from the calling \ac{PE}'s \ac{GPU}; otherwise a
null pointer.
}

\apinotes{
\FUNC{shmemg\_ptr} operates at single-thread granularity. No thread group or
hardware thread group variants are defined. When \VAR{pe} is the calling \ac{PE},
\FUNC{shmemg\_ptr} returns a pointer to the local symmetric data object.
}
\end{apidefinition}

\paragraph*{\textbf{Example}}\mbox{}\par\nobreak
The following program uses \FUNC{shmemg\_ptr} to obtain a directly accessible
pointer to a symmetric buffer on the right neighbor \ac{PE}. When the neighbor
is directly accessible, the kernel writes into the remote buffer with an
ordinary store through the returned pointer; otherwise it falls back to
\FUNC{shmemg\_putmem}.

\lstinputlisting[style=SourceListingsDefault]{example_code/gc_ptr_example.c}

